\chapter{Environnement de développement et méthodologie}

\minitoc

\section{Environnement matériel}
% TODO: décrire brièvement les machines et le matériel utilisés (facultatif).
[Décrivez le matériel utilisé.]

\section{Environnement de travail collaboratif}
% TODO: décrire les outils collaboratifs -- suivi des tâches, gestion des versions, communication.
% Conserver une \subsection par outil. Supprimer cette section pour un projet individuel sans ces outils.

\subsection{[Nom de l'outil, par exemple Trello / Jira]}
[Décrivez l'utilisation de l'outil.]

\subsection{[Nom de l'outil, par exemple GitHub / GitLab]}
[Décrivez l'utilisation de l'outil, par exemple la stratégie de gestion des branches et le processus de revue de code.]

\section{Environnement logiciel}
% TODO: présenter l'environnement de développement intégré, la conteneurisation et les technologies utilisées.

\subsection{[Nom de l'environnement de développement intégré, par exemple Visual Studio Code]}
[Décrivez l'environnement de développement intégré et les extensions utiles.]

\subsection{[Conteneurisation, par exemple Docker]}
[Décrivez son utilisation, le cas échéant.]

\subsection{Technologies}
% TODO: prévoir une \subsubsection par langage, cadriciel, base de données ou bibliothèque utilisés.
% Dupliquer ce bloc pour chaque technologie.

\subsubsection{[Nom de la technologie]}
[Rédigez un court paragraphe présentant la technologie et les raisons de son choix pour ce projet.]

\subsubsection{[Nom de la technologie]}
[Rédigez un court paragraphe.]

\subsection{Conclusion}
% TODO: conclure le chapitre en 2 à 3 phrases.
[Rédigez une brève conclusion du chapitre.]
